% 中文使用 TeX Live 的 Fandol 字体，西文优先使用本机 Times New Roman。
% FandolSong 为宋体字族；无 Times New Roman 时使用 XITS 预览。
\YNUsetup{
  style = {
    font = times,
    cjk-font = fandol,
    font-size = -4,
    hyperlink = none,
    bib-backend = bibtex,
    bib-style = numerical,
    bib-resource = {references.bib}
  }
}

% 本机使用学校指定的西文字体；仅有TeX Live的环境保持原XITS预览。
% 字体文件不随论文源码包分发，来源与文件散列见validation/latin-font-source.json。
\IfFontExistsTF{Times New Roman}{\setmainfont{Times New Roman}}{}

% geometry 的 head/foot 表示页眉高度和页脚基线偏移，不是距纸边距离。
% 保留正文边距；按实际PDF文字边界校准，使页眉上沿、页脚下沿分别
% 距纸边约1.6cm、1.5cm。核验记录见validation/page-layout.json。
\geometry{headsep=0.94cm,footskip=0.92cm}

% 使用模板既有封面接口做两项局部适配，保持上游类文件不变。
% 上游封面忽略 info/date 而直接打印编译年月；此处改为使用配置值。
% Fandol 不含上游封面调用的行楷字族，骨架封面使用同套字体中的楷体。
\ExplSyntaxOn
\EditInstance { YNU / cover } { cover-i-default }
  {
    date / content = \l__YNU_info_date_tl,
    type / format = \zihao { 0 } \use:c { YNU@kai },
    degree / content = \clist_item:Nn \c__YNU_degree_type_clist { \l__YNU_info_degree_type_int }
  }
\ExplSyntaxOff

% 封面标题的人工换行不进入PDF标题元数据。
\AddToHook{package/hyperref/after}{
  \pdfstringdefDisableCommands{\def\\{}}
}
\hypersetup{
  pdfsubject = {硕士学位论文写作稿}
}
